\section{Introduction}

Bluesky and the wider ATProto ecosystem have become increasingly relevant as a source of public social-media data \parencite{failla2024bluesky,quelle2025bluesky}. For smaller languages such as Slovene, this opens up an opportunity to document an emerging online community on a new, decentralised platform, and to consider how such documentation should be done. The paper is also motivated by the author's long-standing interest in Slovene online communities, by familiarity with the emerging Slovene Bluesky space, and by a broader interest in what open-protocol social media makes visible and researchable.

This paper describes an attempt to build a Slovene Bluesky corpus from public data. The main challenges are not only in collection, but in deciding what should count as a Slovene post. Language tags are useful but imperfect, and posts by Slovene-linked users are not automatically Slovene. The approach combines protocol-aware collection with explicit validation, but it still requires drawing a practical line around what counts as a Slovene post.

Four questions guide the work. The first is whether a usable Slovene Bluesky corpus can be built from public ATProto data. The second concerns how far collection can move beyond a simple \texttt{langs=sl} strategy. The third asks which inclusion rules are reliable enough for a main corpus. The fourth asks what basic properties are visible in the resulting corpus.

This paper makes two contributions. First, it presents a protocol-native workflow for constructing a high-precision corpus of Slovene Bluesky posts from public ATProto data, combining tagged discovery, multi-PDS backfill, and validated post-level language filtering. Second, it uses the resulting corpus as an empirical case study of a small online community in formation. The analysis shows a sharp late-2024 growth surge, strongly conversational posting patterns, and a highly concentrated participation structure, together suggesting that the Slovene Bluesky sphere is emerging as a small conversational public space.
